% AI Tools, LLMs & Observability section

% 1 – What is an LLM?
\QuestionSlide[\CategoryBadge[AiToolColor!20]{AI Tools \& LLMs}]{* What is a Large Language Model (LLM)?}

\begin{frame}[fragile]
  \frametitle{\AnswerTitle[\CategoryBadge[AiToolColor!20]{AI Tools \& LLMs}]{What is a Large Language Model?}}

  {\footnotesize
  A \textbf{Large Language Model} is a neural network trained on massive text corpora to predict the next token. At inference, it generates text autoregressively. LLMs exhibit emergent abilities: reasoning, code generation, translation, and tool use. Notable families: GPT-4o (OpenAI), Claude (Anthropic), Gemini (Google), Phi/Mistral (open-weight).
  }
\end{frame}

% 2 – Tokens
\QuestionSlide[\CategoryBadge[AiToolColor!20]{AI Tools \& LLMs}]{* What are tokens in LLMs and why do they matter?}

\begin{frame}[fragile]
  \frametitle{\AnswerTitle[\CategoryBadge[AiToolColor!20]{AI Tools \& LLMs}]{What are tokens in LLMs?}}

  {\footnotesize
  LLMs process text as \textbf{tokens} — sub-word pieces (e.g.\ \texttt{"programming"} $\to$ \texttt{["program","ming"]}). Pricing, context window limits, and latency are all measured in tokens. Rules of thumb: 1 token $\approx$ 0.75 English words; 1 page $\approx$ 500 tokens. The \textbf{context window} is the maximum total tokens (input + output) the model can process at once.
  }
\end{frame}

% 3 – Fine-Tuning vs RAG
\QuestionSlide[\CategoryBadge[AiToolColor!20]{AI Tools \& LLMs}]{** When do you fine-tune vs use RAG?}

\begin{frame}[fragile]
  \frametitle{\AnswerTitle[\CategoryBadge[AiToolColor!20]{AI Tools \& LLMs}]{Fine-tuning vs RAG?}}

  {\footnotesize
  \begin{tabular}{ll}
    \textbf{Fine-Tuning} & \textbf{RAG} \\
    Bakes knowledge into weights & Retrieves context at query time \\
    Expensive, requires labelled data & Cheap, no GPU training \\
    Good for style/format/task & Good for dynamic/private knowledge \\
    Knowledge frozen at train time & Always up-to-date \\
    Hard to audit/update & Easy to add/remove sources \\
  \end{tabular}
  \vspace{0.5em}
  \textbf{Rule}: Use RAG first; fine-tune only for style/domain adaptation when RAG quality is insufficient.
  }
\end{frame}

% 4 – Tool Use / Function Calling
\QuestionSlide[\CategoryBadge[AiToolColor!20]{AI Tools \& LLMs}]{** What is function calling / tool use in LLMs?}

\begin{frame}[fragile]
  \frametitle{\AnswerTitle[\CategoryBadge[AiToolColor!20]{AI Tools \& LLMs}]{What is function calling/tool use?}}

  {\footnotesize
  \textbf{Tool use} (function calling) lets an LLM signal that it wants to invoke a defined function with structured arguments. Your code executes the function, returns the result, and continues the conversation. This gives LLMs real-world agency: database queries, HTTP calls, code execution, sending emails. In .NET: Semantic Kernel Plugins, \texttt{UseFunctionInvocation()} in Microsoft.Extensions.AI.
  }

  \begin{minted}[fontsize=\scriptsize]{csharp}
[KernelFunction, Description("Get current stock price for a ticker symbol")]
public async Task<string> GetStockPrice(
    [Description("Stock ticker, e.g. MSFT")] string ticker)
{
    var price = await _stockApi.GetPriceAsync(ticker);
    return $"{ticker}: {price:C}";
}
  \end{minted}
\end{frame}

% 5 – GitHub Copilot
\QuestionSlide[\CategoryBadge[AiToolColor!20]{AI Tools \& LLMs}]{* What is GitHub Copilot and how does it work?}

\begin{frame}[fragile]
  \frametitle{\AnswerTitle[\CategoryBadge[AiToolColor!20]{AI Tools \& LLMs}]{What is GitHub Copilot?}}

  {\footnotesize
  \textbf{GitHub Copilot} is an AI pair programmer powered by OpenAI Codex / GPT-4o, integrated into VS Code, Visual Studio, JetBrains, and Neovim. It provides \textbf{inline completions}, \textbf{chat} (explain, refactor, fix bugs), \textbf{PR summaries}, and \textbf{Workspace agent} for repo-wide queries. \textbf{Copilot CLI} assists with shell commands.
  }
\end{frame}

% 6 – Claude
\QuestionSlide[\CategoryBadge[AiToolColor!20]{AI Tools \& LLMs}]{* What is Claude and its model family?}

\begin{frame}[fragile]
  \frametitle{\AnswerTitle[\CategoryBadge[AiToolColor!20]{AI Tools \& LLMs}]{What is Claude?}}

  {\footnotesize
  \textbf{Claude} is Anthropic's family of AI assistants designed for safety, instruction-following, and long-context reasoning. Tiers:
  \begin{itemize}
    \item \textbf{Haiku} – fastest and cheapest; ideal for high-volume tasks.
    \item \textbf{Sonnet} – balanced capability and speed; default for most workloads.
    \item \textbf{Opus} – most capable; complex reasoning and agentic tasks.
  \end{itemize}
  Claude supports 200K-token context windows, vision, tool use, and is available via Anthropic API and Azure AI Foundry.
  }
\end{frame}

% 7 – LangChain
\QuestionSlide[\CategoryBadge[AiToolColor!20]{AI Tools \& LLMs}]{** What is LangChain?}

\begin{frame}[fragile]
  \frametitle{\AnswerTitle[\CategoryBadge[AiToolColor!20]{AI Tools \& LLMs}]{What is LangChain?}}

  {\footnotesize
  \textbf{LangChain} (Python / JS) is an open-source framework for composing LLM-powered applications: \textbf{chains} (pipelines of LLM calls), \textbf{agents} (LLMs that select and invoke tools dynamically), \textbf{memory} (conversation history), and \textbf{retrieval} (document loaders, text splitters, vector stores). \textbf{LangGraph} extends it with stateful, graph-based agentic workflows.
  }
\end{frame}

% 8 – Multi-Agent Systems
\QuestionSlide[\CategoryBadge[AiToolColor!20]{AI Tools \& LLMs}]{*** What is a Multi-Agent System in the context of LLMs?}

\begin{frame}[fragile]
  \frametitle{\AnswerTitle[\CategoryBadge[AiToolColor!20]{AI Tools \& LLMs}]{What is a Multi-Agent System?}}

  {\footnotesize
  A \textbf{Multi-Agent System} coordinates multiple specialised AI agents to complete complex tasks. Each agent has a role (Planner, Coder, Reviewer, Researcher) and communicates via messages. Patterns: \textbf{Orchestrator-Worker} (one agent delegates), \textbf{Group Chat} (agents debate), \textbf{Hierarchical} (nested orchestrators). Frameworks: Semantic Kernel \texttt{AgentGroupChat}, AutoGen, LangGraph.
  }

  \begin{minted}[fontsize=\scriptsize]{csharp}
AgentGroupChat groupChat = new(plannerAgent, coderAgent, reviewerAgent)
{
    ExecutionSettings = new()
    {
        SelectionStrategy  = new RoundRobinSelectionStrategy(),
        TerminationStrategy = new ApprovalTerminationStrategy()
    }
};

await foreach (var msg in groupChat.InvokeAsync())
    Console.WriteLine($"[{msg.AuthorName}] {msg.Content}");
  \end{minted}
\end{frame}

% 9 – OpenTelemetry for LLM Observability
\QuestionSlide[\CategoryBadge[AiToolColor!20]{AI Tools \& LLMs}]{*** What is OpenTelemetry-based LLM observability?}

\begin{frame}[fragile]
  \frametitle{\AnswerTitle[\CategoryBadge[AiToolColor!20]{AI Tools \& LLMs}]{OpenTelemetry-based LLM observability?}}

  {\footnotesize
  \textbf{OpenTelemetry for LLMs} captures spans for each model call with standardised attributes: model name, input/output token counts, latency, prompt/completion content (opt-in), and tool calls. The \textbf{GenAI Semantic Conventions} spec (CNCF) defines attribute names like \texttt{gen\_ai.system}, \texttt{gen\_ai.usage.input\_tokens}. \texttt{UseOpenTelemetry()} middleware on \texttt{IChatClient} wires this automatically.
  }

  \begin{minted}[fontsize=\scriptsize]{csharp}
IChatClient client = new OpenAIClient(key)
    .AsChatClient("gpt-4o")
    .AsBuilder()
    .UseOpenTelemetry(
        loggerFactory: loggerFactory,
        configure: t => t.EnableSensitiveData = false) // excludes prompts
    .Build();

// Every call now emits a span to your OTLP collector (Jaeger, Grafana, Azure Monitor)
  \end{minted}
\end{frame}
